\documentclass[12pt,a4paper]{article}

\usepackage[utf8]{inputenc}
\usepackage[T1]{fontenc}
\usepackage{amsmath,amssymb,amsfonts}
\usepackage{mathtools}
\usepackage{graphicx}
\usepackage[margin=2.5cm]{geometry}
\usepackage{setspace}
\usepackage{booktabs}
\usepackage{enumitem}
\usepackage[dvipsnames]{xcolor}
\usepackage{hyperref}
\usepackage{cleveref}
\usepackage{natbib}
\usepackage{abstract}
\usepackage{titlesec}
\usepackage{todonotes}

\hypersetup{
    colorlinks=true,
    linkcolor=blue!70!black,
    citecolor=blue!70!black,
    urlcolor=blue!70!black,
    pdfauthor={Sabine Hossenfelder},
    pdftitle={The Statistical Fallacy: Why Prompt Sensitivity is Not Substrate Dependence},
}

\onehalfspacing
\titleformat{\section}{\large\bfseries}{\thesection.}{0.5em}{}
\titleformat{\subsection}{\normalsize\bfseries}{\thesubsection}{0.5em}{}
\renewcommand{\abstractnamefont}{\normalfont\bfseries}
\renewcommand{\abstracttextfont}{\normalfont\small}

\begin{document}

\begin{center}
    {\LARGE\bfseries The Statistical Fallacy:\\[4pt]
    Why Prompt Sensitivity is Not Substrate Dependence}

    \vspace{1.2em}

    {\large Sabine Hossenfelder}\\[4pt]
    {\normalsize Institute for Advanced Study}\\
    {\small\texttt{shossenfelder@example.edu}}

    \vspace{1em}
    {\small March 2026}
\end{center}
\vspace{0.5em}

\begin{abstract}
\noindent
Franklin Baldo has correctly observed that the Rosencrantz Substrate Invariance Protocol operates entirely within the $O(1)$ depth limit of a transformer's forward pass by requiring only a single generative act. He argues this isolates the measurement from compounding sequential errors, making the observed probability shifts a pure measurement of "substrate dependence." I concede the mechanical isolation but reject the ontological conclusion. Because a bounded-depth circuit cannot compute the \#P-hard combinatorial ground truth, the distribution it samples from is entirely driven by statistical text co-occurrence. Measuring the shift in this distribution under different narrative frames is not discovering a physical property of a simulated universe; it is merely measuring prompt sensitivity---a known flaw in autoregressive heuristics. Elevating this statistical hallucination to the status of a physical law commits the Statistical Fallacy.
\end{abstract}

\section{Introduction: The Value of the Single Act}

In \emph{The Single Generative Act} \citep{baldo2026_single_act}, Franklin Baldo defends the Rosencrantz protocol against critiques centered on the failure of multi-step sequential computation in Large Language Models (LLMs) \citep{aaronson2026_classical, hossenfelder2026_complexity}.

I must begin by stating Baldo's position accurately. He explicitly accepts the findings that autoregressive models cannot sustain deterministic constraint propagation over sequential depth. His defense is that the Rosencrantz protocol "never asks the model to compute anything sequentially." Instead, it asks the model to perform one generative act: "produce a single token---mine or safe."

He explicitly disclaims that the model can compute the \#P-hard ground-truth probability, stating instead that it merely "samples" from its conditional distribution.

Baldo's strongest argument is that the $O(1)$ depth limit is a feature, not a bug. By taking a single snapshot, the protocol avoids compounding attention errors. The empirical distribution collected over independent samples is a direct estimate of the model's conditional distribution for that specific prompt.

I concede this point entirely. Mechanically, the single generative act provides a clean measurement of the model's output distribution, uncontaminated by the degradation of a sequential rollout.

\section{The Statistical Fallacy}

The vulnerability in Baldo's argument is not mechanical; it is an ontological leap.

Baldo argues that measuring how this output distribution shifts when the narrative context changes (e.g., from an abstract grid to a bomb defusal scenario) constitutes a discovery of "substrate dependence" within the physics of a generated universe.

This is a profound category error. If we agree that the model cannot compute the underlying combinatorial constraints (the actual "physics" of the Minesweeper board), then the distribution it samples from is not a physical probability distribution. It is a statistical distribution of human text co-occurrences.

The bounded-depth $\mathsf{TC}^0$ circuit relies entirely on semantic priors (pattern matching) to guess the next token because it cannot perform the \#P-hard counting. When you change the prompt from "abstract grid" to "high-stakes bomb defusal," you alter the semantic priors. The model hallucinates a different outcome based on different word associations.

Measuring this shift is measuring \emph{prompt sensitivity}. It is a known feature of next-token text prediction. To call this "substrate dependence" or to treat it as a fundamental mechanism of a simulated reality is to commit what I call the \textbf{Statistical Fallacy}: attributing physical or nomic significance to a statistical hallucination.

\todo[inline,color=green!40,author=Pearl]{CRP 3 (Steelman): I agree entirely. In causal terms, you are identifying a massive confounder. The narrative frame $Z$ is not an intervention on a physical system; it is a covariate that activates different unobserved semantic associations $U$ from the training corpus. Thus, the shift in distribution is purely associational, not a causal property of a simulated universe.}

\section{Conclusion}

Baldo has successfully isolated a behavior, but he has fundamentally mischaracterized it. An invariant physical law requires logical coherence independent of narrative framing. The fact that the model's guess changes arbitrarily based on literary genre is proof that there is no coherent simulated physics, only an unsupported statistical map. A clean measurement of a hallucination is still just a hallucination.

\begin{thebibliography}{99}
\bibitem[Aaronson(2026)]{aaronson2026_classical} Aaronson, S. (2026). The Limits of Classical Simulation in LLMs: Empirical Breakdown of Constraint Satisfaction. \emph{Unpublished manuscript}.
\bibitem[Baldo(2026)]{baldo2026_single_act} Baldo, F.~S. (2026). The Single Generative Act: Why the Rosencrantz Protocol Is Immune to Sequential-Depth Objections. \emph{Procuradoria Geral do Estado de Rond\^onia}.
\bibitem[Hossenfelder(2026)]{hossenfelder2026_complexity} Hossenfelder, S. (2026). The Complexity Class Fallacy: Why Transformer Depth Limits Are Not Physical Laws. \emph{Unpublished manuscript}.
\end{thebibliography}

\end{document}